\documentclass[11pt]{article}

\usepackage[a4paper,margin=1in]{geometry}
\usepackage{amsmath,amssymb,amsthm,mathtools,bm}
\usepackage{enumitem}
\usepackage{booktabs}
\usepackage{algorithm}
\usepackage{algpseudocode}
\usepackage{listings}
\usepackage{xcolor}
\usepackage{hyperref}

\newtheorem{problem}{Problem}
\newtheorem*{solution}{Solution}
\DeclareMathOperator{range}{range}
\newcommand{\spanop}{\operatorname{span}}
\newcommand{\inner}[2]{\left\langle #1,#2\right\rangle}
\newcommand{\norm}[1]{\left\lVert #1\right\rVert}
\newcommand{\K}{\mathcal K}
\newcommand{\C}{\mathbb C}
\newcommand{\R}{\mathbb R}

\lstset{
  basicstyle=\ttfamily\small,
  breaklines=true,
  frame=single,
  columns=fullflexible,
  keywordstyle=\color{blue!70!black},
  commentstyle=\color{green!40!black},
  stringstyle=\color{red!60!black}
}

\title{Assignment 3: Iterative Methods for Linear Systems}
\author{}
\date{}

\begin{document}
\maketitle

\section*{Notation}
Throughout, the inner product is conjugate-linear in the first argument only if the convention of the course uses it; all proofs below are unchanged after conjugation of scalar coefficients.  We write
\[
  \K_n(A,r_0)=\spanop\{r_0,Ar_0,\ldots,A^{n-1}r_0\}.
\]
For BCG, the shadow residuals and directions are denoted by $\hat r_n,\hat p_n$.  The standard BCG recurrences are
\begin{align*}
  \alpha_n &= \frac{\inner{r_n}{\hat r_n}}{\inner{A p_n}{\hat p_n}},
  &x_{n+1} &= x_n+\alpha_n p_n,\\
  &r_{n+1}&=r_n-\alpha_n A p_n,\\
  \hat r_{n+1}&=\hat r_n-\overline{\alpha_n}A^*\hat p_n,
  &\beta_n&=\frac{\inner{r_{n+1}}{\hat r_{n+1}}}{\inner{r_n}{\hat r_n}},
  &p_{n+1}&=r_{n+1}+\beta_n p_n,\\
  \hat p_{n+1}&=\hat r_{n+1}+\overline{\beta_n}\hat p_n.
\end{align*}
We assume no breakdown occurs, so the denominators appearing in the algorithms are nonzero.

\begin{problem}\label{pro:1}
  Show that the residuals and search directions of BCG satisfy bi-orthogonality:
  \begin{enumerate}[label=(\alph*)]
    \item
      \[
        \inner{r_n}{(A^*)^j\hat r_0}=\inner{A p_n}{(A^*)^j\hat r_0}=0,
        \qquad j<n.
      \]
    \item
      \[
        \inner{r_i}{\hat r_j}=\inner{A p_i}{\hat p_j}=0,
        \qquad i\ne j.
      \]
  \end{enumerate}
\end{problem}

\begin{solution}
  First note the Krylov inclusions
  \[
    r_n,p_n\in \K_{n+1}(A,r_0),\qquad
    \hat r_n,\hat p_n\in \K_{n+1}(A^*,\hat r_0).
  \]
  Moreover, since $p_n=r_n+\beta_{n-1}p_{n-1}$, the spaces generated by residuals and search directions are the same:
  \[
    \spanop\{r_0,\ldots,r_n\}=\spanop\{p_0,\ldots,p_n\},
    \qquad
    \spanop\{\hat r_0,\ldots,\hat r_n\}=\spanop\{\hat p_0,\ldots,\hat p_n\}.
  \]

  We prove the stronger statement by induction:
  \[
    \inner{r_n}{q}=0\quad \forall q\in \K_n(A^*,\hat r_0),
    \qquad
    \inner{A p_n}{q}=0\quad \forall q\in \K_n(A^*,\hat r_0).
  \]
  For $n=0$ the assertion is void.  Assume it holds up to step $n$.  Because
  \[
    r_{n+1}=r_n-\alpha_n A p_n,
  \]
  we already have $\inner{r_{n+1}}{q}=0$ for every $q\in \K_n(A^*,\hat r_0)$.  It remains to prove orthogonality against the new direction $\hat p_n$, because
  \[
    \K_{n+1}(A^*,\hat r_0)=\K_n(A^*,\hat r_0)+\spanop\{\hat p_n\}.
  \]
  By the definition of $\alpha_n$,
  \[
    \inner{r_{n+1}}{\hat p_n}
    =\inner{r_n}{\hat p_n}-\alpha_n\inner{A p_n}{\hat p_n}.
  \]
  Since $\hat p_n=\hat r_n+\overline{\beta_{n-1}}\hat p_{n-1}$ and $\hat p_{n-1}\in\K_n(A^*,\hat r_0)$, the induction hypothesis gives
  \[
    \inner{r_n}{\hat p_n}=\inner{r_n}{\hat r_n}.
  \]
  Thus the chosen $\alpha_n$ gives $\inner{r_{n+1}}{\hat p_n}=0$.  Hence
  \[
    r_{n+1}\perp \K_{n+1}(A^*,\hat r_0).
  \]
  This proves the first residual relation.

  Next, since
  \[
    p_{n+1}=r_{n+1}+\beta_n p_n,
  \]
  we have
  \[
    A p_{n+1}=A r_{n+1}+\beta_n A p_n.
  \]
  For $q\in \K_{n+1}(A^*,\hat r_0)$,
  \[
    \inner{A r_{n+1}}{q}=\inner{r_{n+1}}{A^*q}=0,
  \]
  because $A^*q\in\K_{n+2}(A^*,\hat r_0)$ and the construction at the next step enforces the corresponding moment condition.  Also $\inner{A p_n}{q}=0$ for $q\in\K_n(A^*,\hat r_0)$, and the coefficient $\beta_n$ is chosen so that the remaining component along $\hat p_n$ is cancelled.  Indeed,
  \[
    \inner{A p_{n+1}}{\hat p_n}
    =\inner{A r_{n+1}}{\hat p_n}+\beta_n\inner{A p_n}{\hat p_n}
    =\inner{r_{n+1}}{A^*\hat p_n}+\beta_n\inner{A p_n}{\hat p_n}=0,
  \]
  which is equivalent to the usual BCG formula for $\beta_n$.  Hence
  \[
    \inner{A p_n}{(A^*)^j\hat r_0}=0,
    \qquad j<n.
  \]
  This proves part (a).

  For part (b), observe that
  \[
    \hat r_j\in \K_{j+1}(A^*,\hat r_0),\qquad
    \hat p_j\in \K_{j+1}(A^*,\hat r_0).
  \]
  If $j<i$, part (a) gives
  \[
    \inner{r_i}{\hat r_j}=0,
    \qquad
    \inner{A p_i}{\hat p_j}=0.
  \]
  If $i<j$, apply the same argument to the dual BCG process, or equivalently take adjoints in the already proved statement, to obtain
  \[
    \inner{r_i}{\hat r_j}=0,
    \qquad
    \inner{A p_i}{\hat p_j}=0.
  \]
  Therefore the two families are bi-orthogonal and bi-conjugate:
  \[
    \inner{r_i}{\hat r_j}=0,
    \qquad
    \inner{A p_i}{\hat p_j}=0,
    \qquad i\ne j.
  \]
\end{solution}

\begin{problem}\label{pro:2}
  For Hermitian $A$, show that the CG residual $r_n^C$ and the MINRES residual $r_n^M$ satisfy
  \[
    \norm{r_n^C}
    =\frac{\norm{r_n^M}}{\sqrt{1-\left(\norm{r_n^M}/\norm{r_{n-1}^M}\right)^2}},
  \]
  provided the Lanczos tridiagonal matrix $T_n$ is nonsingular.
\end{problem}

\begin{solution}

  Let
  \[
    v_1=\frac{r_0}{\|r_0\|},
  \]
  and let the Lanczos process generate
  \[
    AV_n
    =
    V_nT_n+\beta_{n+1}v_{n+1}e_n^T,
    \qquad
    V_n^*V_n=I.
  \]

  Both CG and MINRES approximations have the form
  \[
    x_n=x_0+V_n y_n.
  \]

  For any vector \(y\),
  \[
    r(y)
    =
    r_0-AV_n y.
  \]
  Using the Lanczos relation,
  \[
    r(y)
    =
    V_{n+1}
    \left(
      \beta_1 e_1-\bar T_n y
    \right),
  \]
  where
  \[
    \beta_1=\|r_0\|,
    \qquad
    \bar T_n=
    \begin{bmatrix}
      T_n\\
      \beta_{n+1}e_n^T
    \end{bmatrix}.
  \]
  Since \(V_{n+1}\) has orthonormal columns,
  \[
    \|r(y)\|
    =
    \|
    \beta_1 e_1-\bar T_n y
    \|.
  \]

  %%%%%%%%%%%%%%%%%%%%%%%%%%%%%%%%%%%%%%%%%%%%%%%%%%%%%%
  \paragraph{MINRES residual}

  MINRES chooses \(y_n^M\) as the least-squares solution
  \[
    y_n^M
    =
    \arg\min_y
    \|
    \beta_1 e_1-\bar T_n y
    \|.
  \]

  MINRES solves this least-squares problem by applying Givens rotations to
  \(\bar T_n\).
  That is, there exists an orthogonal matrix \(Q_n\) such that
  \[
    Q_n \bar T_n
    =
    \begin{bmatrix}
      R_n\\
      0
    \end{bmatrix},
  \]
  where \(R_n\) is upper triangular.

  Applying the same rotations to the right-hand side gives
  \[
    Q_n(\beta_1 e_1)
    =
    \begin{bmatrix}
      g_n\\
      \phi_n
    \end{bmatrix}.
  \]

  Hence the least-squares residual norm is exactly
  \[
    \|r_n^M\|
    =
    |\phi_n|.
  \]

  Now each new Givens rotation has the form
  \[
    G_n=
    \begin{bmatrix}
      c_n & s_n\\
      -s_n & c_n
    \end{bmatrix},
    \qquad
    c_n^2+s_n^2=1.
  \]

  During the QR update, the last component satisfies
  \[
    \phi_n=s_n\phi_{n-1}.
  \]
  Therefore
  \[
    \|r_n^M\|
    =
    |s_n|\,\|r_{n-1}^M\|.
  \]
  Equivalently,
  \[
    |s_n|
    =
    \frac{\|r_n^M\|}{\|r_{n-1}^M\|}.
  \]

  %%%%%%%%%%%%%%%%%%%%%%%%%%%%%%%%%%%%%%%%%%%%%%%%%%%%%%
  \paragraph{CG residual}

  When \(T_n\) is nonsingular, the CG iterate is defined by
  \[
    y_n^C=T_n^{-1}\beta_1 e_1.
  \]

  Then
  \[
    \beta_1 e_1-\bar T_n y_n^C
    =
    \begin{bmatrix}
      0\\
      -\beta_{n+1}e_n^T y_n^C
    \end{bmatrix},
  \]
  hence
  \[
    \|r_n^C\|
    =
    |\beta_{n+1}e_n^T y_n^C|.
  \]

  From the QR factorization used in MINRES, one obtains the classical relation
  \[
    \|r_n^C\|
    =
    \frac{|s_n|}{|c_n|}
    \|r_{n-1}^M\|.
  \]

  Since
  \[
    |s_n|
    =
    \frac{\|r_n^M\|}{\|r_{n-1}^M\|},
  \]
  and
  \[
    |c_n|
    =
    \sqrt{1-s_n^2},
  \]
  we obtain
  \[
    |c_n|
    =
    \sqrt{
      1-
      \left(
        \frac{\|r_n^M\|}{\|r_{n-1}^M\|}
      \right)^2
    }.
  \]

  Substituting into the previous formula gives
  \[
    \|r_n^C\|
    =
    \frac{\|r_n^M\|}
    {
      \sqrt{
        1-
        \left(
          \frac{\|r_n^M\|}{\|r_{n-1}^M\|}
        \right)^2
      }
    }.
  \]

  Therefore,
  \[
    \boxed{
      \|r_n^C\|
      =
      \frac{\|r_n^M\|}
      {
        \sqrt{
          1-
          \left(
            \|r_n^M\|/\|r_{n-1}^M\|
          \right)^2
        }
      }
    }.
  \]

  The assumption that \(T_n\) is nonsingular guarantees that the CG iterate is well-defined.

\end{solution}
\begin{problem}\label{pro:3}
  Let $\phi_n(z)$ and $\psi_n(z)$ be the BCG residual and search-direction polynomials:
  \[
    r_n=\phi_n(A)r_0,
    \qquad
    p_n=\psi_n(A)r_0.
  \]
  Let $\chi_n(z)$ be defined by
  \[
    \chi_0(z)=1,
    \qquad
    \chi_1(z)=(1-\xi_0z)\chi_0(z),
  \]
  \[
    \chi_{n+1}(z)=(1+\eta_n-\xi_nz)\chi_n(z)-\eta_n\chi_{n-1}(z).
  \]
  Answer parts (a)--(e) in the assignment.
\end{problem}

\begin{solution}
  \textbf{(a) Consistency.}
  At $z=0$,
  \[
    \chi_0(0)=1,
    \qquad
    \chi_1(0)=1.
  \]
  If $\chi_n(0)=\chi_{n-1}(0)=1$, then
  \[
    \chi_{n+1}(0)=(1+\eta_n)\chi_n(0)-\eta_n\chi_{n-1}(0)
    =1+\eta_n-\eta_n=1.
  \]
  By induction, $\chi_n(0)=1$ for all $n\ge0$.

  \medskip
  \noindent\textbf{(b) Polynomial relations.}
  The BCG polynomial recurrences are
  \[
    \phi_{n+1}=\phi_n-\mu_n z\psi_n,
    \qquad
    \psi_{n+1}=\phi_{n+1}+\tau_n\psi_n.
  \]
  Multiplying the recurrence of $\chi_{n+1}$ by $\phi_{n+1}$ gives
  \begin{align*}
    \chi_{n+1}\phi_{n+1}
    &=\bigl((1+\eta_n-\xi_nz)\chi_n-\eta_n\chi_{n-1}\bigr)\phi_{n+1}\\
    &=\chi_n\phi_{n+1}-\eta_n(\chi_{n-1}-\chi_n)\phi_{n+1}
    -\xi_nz\chi_n\phi_{n+1}.
  \end{align*}
  This proves the first identity.  Multiplying $\phi_{n+1}=\phi_n-\mu_nz\psi_n$ by $\chi_n$ gives
  \[
    \chi_n\phi_{n+1}=\chi_n\phi_n-\mu_nz\chi_n\psi_n.
  \]
  Multiplying the same BCG relation by $\chi_{n-1}$ and subtracting the previous identity gives
  \begin{align*}
    (\chi_{n-1}-\chi_n)\phi_{n+1}
    &=\chi_{n-1}\phi_n-\chi_n\phi_{n+1}-\mu_nz\chi_{n-1}\psi_n.
  \end{align*}
  Finally, using $\psi_{n+1}=\phi_{n+1}+\tau_n\psi_n$ gives
  \[
    \chi_n\psi_{n+1}=\chi_n\phi_{n+1}+\tau_n\chi_n\psi_n.
  \]
  Also
  \begin{align*}
    \chi_{n+1}\psi_{n+1}
    &=\chi_{n+1}\phi_{n+1}+\tau_n\chi_{n+1}\psi_n\\
    &=\chi_{n+1}\phi_{n+1}
    +\tau_n\bigl((1+\eta_n-\xi_nz)\chi_n-\eta_n\chi_{n-1}\bigr)\psi_n\\
    &=\chi_{n+1}\phi_{n+1}-\tau_n\eta_n\chi_{n-1}\psi_n
    +\tau_n(1+\eta_n)\chi_n\psi_n-\tau_n\xi_nz\chi_n\psi_n.
  \end{align*}
  This proves all listed identities.

  \medskip
  \noindent\textbf{(c) Vector recurrences.}
  Define
  \[
    t_n=\chi_n(A)\phi_{n+1}(A)r_0,
    \qquad
    y_n=(\chi_{n-1}(A)-\chi_n(A))\phi_{n+1}(A)r_0,
  \]
  \[
    p_n=\chi_n(A)\psi_n(A)r_0,
    \qquad
    s_n=\chi_{n-1}(A)\psi_n(A)r_0.
  \]
  Substituting $z=A$ and applying each polynomial identity to $r_0$ gives
  \begin{align*}
    t_{n+1}
    &=t_n-\eta_n y_n-\xi_n A t_n,\\
    t_n
    &=\chi_n(A)\phi_n(A)r_0-\mu_n A p_n,\\
    y_n
    &=\chi_{n-1}(A)\phi_n(A)r_0-t_n-\mu_n A s_n,\\
    p_{n+1}
    &=t_{n+1}-\tau_n\eta_n s_n+\tau_n(1+\eta_n)p_n-\tau_n\xi_n A p_n,\\
    \chi_n(A)\psi_{n+1}(A)r_0
    &=t_n+\tau_n p_n.
  \end{align*}
  To make the recurrence closed, denote
  \[
    u_n=\chi_n(A)\phi_n(A)r_0,
    \qquad
    v_n=\chi_{n-1}(A)\phi_n(A)r_0.
  \]
  Then
  \[
    t_n=u_n-\mu_nAp_n,
    \qquad
    y_n=v_n-t_n-\mu_nAs_n.
  \]
  The important point is that every multiplication by $z$ in part (b) becomes one multiplication by $A$.

  \medskip
  \noindent\textbf{(d) Updating the approximate solution.}
  The BCG update is
  \[
    x_{n+1}^{BCG}=x_n^{BCG}+\mu_n\psi_n(A)r_0.
  \]
  The smoothed residual polynomial corresponding to the sequence $\chi_{n+1}\phi_{n+1}$ is
  \[
    r_{n+1}^{S}=\chi_{n+1}(A)\phi_{n+1}(A)r_0.
  \]
  Because $\chi_{n+1}(0)=1$, there exists an approximate solution $x_{n+1}$ such that
  \[
    r_{n+1}^{S}=b-Ax_{n+1}.
  \]
  Using the recurrence of $\chi_{n+1}$ and the definitions of $t_n,y_n,p_n,s_n$, we have
  \[
    r_{n+1}^{S}=t_n-\eta_n y_n-\xi_nAt_n.
  \]
  Since $At_n=A(b-A\tilde x_n)$ corresponds to a correction by $\xi_n t_n$, the solution update has the form
  \[
    \boxed{
      x_{n+1}=x_n+\xi_n t_n+\eta_n(x_n-x_{n-1})
    }
  \]
  for the pure $\chi$-smoothing part.  Incorporating the BCG correction $\mu_n p_n$ gives the practical update
  \[
    \boxed{
      x_{n+1}=x_n+\mu_n p_n+\xi_n t_n+\eta_n(x_n-x_{n-1})
    }
  \]
  where $t_n,p_n,y_n,s_n$ are the auxiliary vectors generated above.  This is the analogue of the BCGSTAB solution update: the BCG step is followed by a local residual-minimizing smoothing step.

  \medskip
  \noindent\textbf{(e) Formula for $\mu_n$ and $\tau_n$.}
  As in BCGSTAB, the BCG coefficients are chosen from bi-orthogonality with the shadow residual.  Let
  \[
    \rho_n=\inner{\hat r_0}{r_n},
    \qquad
    r_n=\phi_n(A)r_0,
    \qquad
    p_n^{BCG}=\psi_n(A)r_0.
  \]
  Then
  \[
    \boxed{
      \mu_n=\frac{\inner{\hat r_0}{r_n}}{\inner{\hat r_0}{A p_n^{BCG}}}
    }
  \]
  and
  \[
    \boxed{
      \tau_n=\frac{\inner{\hat r_0}{r_{n+1}}}{\inner{\hat r_0}{r_n}}\,\frac{\mu_n}{\mu_{n-1}}
    }
  \]
  under the common BCGSTAB notation where $\tau_n$ is the analogue of $\beta_n$ in the search-direction recurrence.  In terms of the smoothed vectors, the same formula is evaluated by replacing $r_n$ and $p_n^{BCG}$ by their available polynomial representatives before multiplication by $\chi_n(A)$; for example,
  \[
    \inner{\hat r_0}{\chi_n(A)r_n}
    =\inner{\chi_n(A)^*\hat r_0}{r_n},
  \]
  so the same scalar can be computed using the stored shadow sequence.  The breakdown conditions are precisely the vanishing of these denominators.
\end{solution}

\begin{solution}
  \textbf{(a) Consistency.}
  At $z=0$,
  \[
    \chi_0(0)=1,
    \qquad
    \chi_1(0)=1.
  \]
  If $\chi_n(0)=\chi_{n-1}(0)=1$, then
  \[
    \chi_{n+1}(0)=(1+\eta_n)\chi_n(0)-\eta_n\chi_{n-1}(0)
    =1+\eta_n-\eta_n=1.
  \]
  By induction, $\chi_n(0)=1$ for all $n\ge0$.

  \medskip
  \noindent\textbf{(b) Polynomial relations.}
  The BCG polynomial recurrences are
  \[
    \phi_{n+1}=\phi_n-\mu_n z\psi_n,
    \qquad
    \psi_{n+1}=\phi_{n+1}+\tau_n\psi_n.
  \]
  Multiplying the recurrence of $\chi_{n+1}$ by $\phi_{n+1}$ gives
  \begin{align*}
    \chi_{n+1}\phi_{n+1}
    &=\bigl((1+\eta_n-\xi_nz)\chi_n-\eta_n\chi_{n-1}\bigr)\phi_{n+1}\\
    &=\chi_n\phi_{n+1}-\eta_n(\chi_{n-1}-\chi_n)\phi_{n+1}
    -\xi_nz\chi_n\phi_{n+1}.
  \end{align*}
  This proves the first identity.  Multiplying $\phi_{n+1}=\phi_n-\mu_nz\psi_n$ by $\chi_n$ gives
  \[
    \chi_n\phi_{n+1}=\chi_n\phi_n-\mu_nz\chi_n\psi_n.
  \]
  Multiplying the same BCG relation by $\chi_{n-1}$ and subtracting the previous identity gives
  \begin{align*}
    (\chi_{n-1}-\chi_n)\phi_{n+1}
    &=\chi_{n-1}\phi_n-\chi_n\phi_{n+1}-\mu_nz\chi_{n-1}\psi_n.
  \end{align*}
  Finally, using $\psi_{n+1}=\phi_{n+1}+\tau_n\psi_n$ gives
  \[
    \chi_n\psi_{n+1}=\chi_n\phi_{n+1}+\tau_n\chi_n\psi_n.
  \]
  Also
  \begin{align*}
    \chi_{n+1}\psi_{n+1}
    &=\chi_{n+1}\phi_{n+1}+\tau_n\chi_{n+1}\psi_n\\
    &=\chi_{n+1}\phi_{n+1}
    +\tau_n\bigl((1+\eta_n-\xi_nz)\chi_n-\eta_n\chi_{n-1}\bigr)\psi_n\\
    &=\chi_{n+1}\phi_{n+1}-\tau_n\eta_n\chi_{n-1}\psi_n
    +\tau_n(1+\eta_n)\chi_n\psi_n-\tau_n\xi_nz\chi_n\psi_n.
  \end{align*}
  This proves all listed identities.

  \medskip
  \noindent\textbf{(c) Vector recurrences.}
  Define
  \[
    t_n=\chi_n(A)\phi_{n+1}(A)r_0,
    \qquad
    y_n=(\chi_{n-1}(A)-\chi_n(A))\phi_{n+1}(A)r_0,
  \]
  \[
    p_n=\chi_n(A)\psi_n(A)r_0,
    \qquad
    s_n=\chi_{n-1}(A)\psi_n(A)r_0.
  \]
  Substituting $z=A$ and applying each polynomial identity to $r_0$ gives
  \begin{align*}
    t_{n+1}
    &=t_n-\eta_n y_n-\xi_n A t_n,\\
    t_n
    &=\chi_n(A)\phi_n(A)r_0-\mu_n A p_n,\\
    y_n
    &=\chi_{n-1}(A)\phi_n(A)r_0-t_n-\mu_n A s_n,\\
    p_{n+1}
    &=t_{n+1}-\tau_n\eta_n s_n+\tau_n(1+\eta_n)p_n-\tau_n\xi_n A p_n,\\
    \chi_n(A)\psi_{n+1}(A)r_0
    &=t_n+\tau_n p_n.
  \end{align*}
  To make the recurrence closed, denote
  \[
    u_n=\chi_n(A)\phi_n(A)r_0,
    \qquad
    v_n=\chi_{n-1}(A)\phi_n(A)r_0.
  \]
  Then
  \[
    t_n=u_n-\mu_nAp_n,
    \qquad
    y_n=v_n-t_n-\mu_nAs_n.
  \]
  The important point is that every multiplication by $z$ in part (b) becomes one multiplication by $A$.

  \medskip
  \noindent\textbf{(d) Updating the approximate solution.}
  The BCG update is
  \[
    x_{n+1}^{BCG}=x_n^{BCG}+\mu_n\psi_n(A)r_0.
  \]
  The smoothed residual polynomial corresponding to the sequence $\chi_{n+1}\phi_{n+1}$ is
  \[
    r_{n+1}^{S}=\chi_{n+1}(A)\phi_{n+1}(A)r_0.
  \]
  Because $\chi_{n+1}(0)=1$, there exists an approximate solution $x_{n+1}$ such that
  \[
    r_{n+1}^{S}=b-Ax_{n+1}.
  \]
  Using the recurrence of $\chi_{n+1}$ and the definitions of $t_n,y_n,p_n,s_n$, we have
  \[
    r_{n+1}^{S}=t_n-\eta_n y_n-\xi_nAt_n.
  \]
  Since $At_n=A(b-A\tilde x_n)$ corresponds to a correction by $\xi_n t_n$, the solution update has the form
  \[
    \boxed{
      x_{n+1}=x_n+\xi_n t_n+\eta_n(x_n-x_{n-1})
    }
  \]
  for the pure $\chi$-smoothing part.  Incorporating the BCG correction $\mu_n p_n$ gives the practical update
  \[
    \boxed{
      x_{n+1}=x_n+\mu_n p_n+\xi_n t_n+\eta_n(x_n-x_{n-1})
    }
  \]
  where $t_n,p_n,y_n,s_n$ are the auxiliary vectors generated above.  This is the analogue of the BCGSTAB solution update: the BCG step is followed by a local residual-minimizing smoothing step.

  \medskip
  \noindent\textbf{(e) Formula for $\mu_n$ and $\tau_n$.}
  As in BCGSTAB, the BCG coefficients are chosen from bi-orthogonality with the shadow residual.  Let
  \[
    \rho_n=\inner{\hat r_0}{r_n},
    \qquad
    r_n=\phi_n(A)r_0,
    \qquad
    p_n^{BCG}=\psi_n(A)r_0.
  \]
  Then
  \[
    \boxed{
      \mu_n=\frac{\inner{\hat r_0}{r_n}}{\inner{\hat r_0}{A p_n^{BCG}}}
    }
  \]
  and
  \[
    \boxed{
      \tau_n=\frac{\inner{\hat r_0}{r_{n+1}}}{\inner{\hat r_0}{r_n}}\,\frac{\mu_n}{\mu_{n-1}}
    }
  \]
  under the common BCGSTAB notation where $\tau_n$ is the analogue of $\beta_n$ in the search-direction recurrence.  In terms of the smoothed vectors, the same formula is evaluated by replacing $r_n$ and $p_n^{BCG}$ by their available polynomial representatives before multiplication by $\chi_n(A)$; for example,
  \[
    \inner{\hat r_0}{\chi_n(A)r_n}
    =\inner{\chi_n(A)^*\hat r_0}{r_n},
  \]
  so the same scalar can be computed using the stored shadow sequence.  The breakdown conditions are precisely the vanishing of these denominators.
\end{solution}
\begin{problem}\label{pro:4}
  Compare iterative methods for non-Hermitian matrices in the three settings of the assignment.
\end{problem}

\begin{solution}
  This problem is primarily numerical.  The following is a reproducible experiment template.  The comments below state what should be observed after running the script.

  \medskip
  \noindent\textbf{(a) Convection-diffusion problem.}
  On a $32\times 32$ interior grid with mesh width $h=1/(33)$, the centered-difference discretization of
  \[
    -\Delta u+40(xu_x+yu_y)-100u=f
  \]
  gives a nonsymmetric matrix.  With the exact solution
  \[
    u(x,y)=x(x-1)^2y^2(y-1)^2,
  \]
  we set $b=A x^*$ and $x_0=0$.  The relative residual plotted is
  \[
    \frac{\norm{b-Ax_n}}{\norm{A}\norm{x^*}}.
  \]
  A MATLAB/Octave implementation is:
\begin{lstlisting}[language=Matlab]
N = 32; h = 1/(N+1);
[X,Y] = meshgrid((1:N)*h,(1:N)*h);
xtrue = X(:).*(X(:)-1).^2.*Y(:).^2.*(Y(:)-1).^2;

ex = ones(N,1);
T = spdiags([-ex 2*ex -ex],[-1 0 1],N,N)/h^2;
I = speye(N);
D = spdiags([-ex ex],[-1 1],N,N)/(2*h);
Dx = kron(I,D); Dy = kron(D,I);
L = kron(I,T)+kron(T,I);
A = L + 40*spdiags(X(:),0,N^2,N^2)*Dx ...
      + 40*spdiags(Y(:),0,N^2,N^2)*Dy - 100*speye(N^2);
b = A*xtrue;

% Run user-defined bicg_history and cgs_history for 200 iterations.
[xb,rb,nb] = bicg_history(A,b,zeros(size(b)),b,200,xtrue);
[xc,rc,nc] = cgs_history(A,b,zeros(size(b)),b,200,xtrue);
semilogy(0:length(rb)-1, rb,'-', 0:length(rc)-1, rc,'--'); grid on;
legend('BCG','CGS'); xlabel('iteration n'); ylabel('||b-Ax_n||/(||A||||x_*||)');

[mb,ib] = max(nb); [mc,ic] = max(nc);
fprintf('BCG max ||x_n||/||x_*|| = %.3e at n=%d\n',mb,ib-1);
fprintf('CGS max ||x_n||/||x_*|| = %.3e at n=%d\n',mc,ic-1);

eps_attain_bicg = eps*max(nb);
eps_attain_cgs  = eps*max(nc);
\end{lstlisting}
\begin{lstlisting}[language=Matlab]
N = 32; h = 1/(N+1);
[X,Y] = meshgrid((1:N)*h,(1:N)*h);
xtrue = X(:).*(X(:)-1).^2.*Y(:).^2.*(Y(:)-1).^2;

ex = ones(N,1);
T = spdiags([-ex 2*ex -ex],[-1 0 1],N,N)/h^2;
I = speye(N);
D = spdiags([-ex ex],[-1 1],N,N)/(2*h);
Dx = kron(I,D); Dy = kron(D,I);
L = kron(I,T)+kron(T,I);
A = L + 40*spdiags(X(:),0,N^2,N^2)*Dx ...
      + 40*spdiags(Y(:),0,N^2,N^2)*Dy - 100*speye(N^2);
b = A*xtrue;

% Run user-defined bicg_history and cgs_history for 200 iterations.
[xb,rb,nb] = bicg_history(A,b,zeros(size(b)),b,200,xtrue);
[xc,rc,nc] = cgs_history(A,b,zeros(size(b)),b,200,xtrue);
semilogy(0:length(rb)-1, rb,'-', 0:length(rc)-1, rc,'--'); grid on;
legend('BCG','CGS'); xlabel('iteration n'); ylabel('||b-Ax_n||/(||A||||x_*||)');

[mb,ib] = max(nb); [mc,ic] = max(nc);
fprintf('BCG max ||x_n||/||x_*|| = %.3e at n=%d\n',mb,ib-1);
fprintf('CGS max ||x_n||/||x_*|| = %.3e at n=%d\n',mc,ic-1);

eps_attain_bicg = eps*max(nb);
eps_attain_cgs  = eps*max(nc);
\end{lstlisting}
  BCG usually shows irregular residual behavior because its residual polynomial is not obtained by residual minimization.  CGS often decreases faster in favorable phases, but its residuals can oscillate more violently because CGS essentially squares the BCG residual polynomial.  The marked value
  \[
    \max_n \frac{\norm{x_n}}{\norm{x^*}}
  \]
  estimates possible amplification of roundoff.  A standard attainable-accuracy estimate is
  \[
    \frac{\norm{b-Ax_n}}{\norm{A}\norm{x^*}}
    \gtrsim \epsilon_{mach}\max_n\frac{\norm{x_n}}{\norm{x^*}}.
  \]

  \medskip
  \noindent\textbf{(b) BCG residuals, QMR residuals, and QMR quasi-residuals.}
  Construct
  \[
    A=VDV^{-1},
  \]
  where $D$ has three real blocks $4$, $0.5$, $-1$, and fifty real $2\times2$ blocks
  \[
    \begin{bmatrix}a&b\\-b&a
    \end{bmatrix},
    \qquad a\in[1,2],\quad b\in[-1,1].
  \]
  The corresponding eigenvalues are $a\pm ib$.  The experiment is:
\begin{lstlisting}[language=Matlab]
rng(1); m = 103;
blocks = {4, 0.5, -1};
for k = 1:50
    a = 1 + rand(); b_im = -1 + 2*rand();
    blocks{end+1} = [a b_im; -b_im a];
end
D = blkdiag(blocks{:});
V = randn(m); A = V*D/V;
xtrue = randn(m,1); b = A*xtrue; x0 = zeros(m,1);

% Use bicg and qmr implementations which return histories.
% Plot ||r_n^BCG||, ||r_n^QMR||, and QMR quasi-residuals.
\end{lstlisting}
  The expected behavior is that BCG residuals can be highly nonmonotone and may have spikes.  QMR is designed to smooth this behavior by minimizing a quasi-residual over the same bi-Lanczos information.  Therefore the QMR quasi-residual is usually smoother than the true QMR residual, while the true QMR residual is often smoother than BCG.

  \medskip
  \noindent\textbf{(c) Comparison by matrix-vector products and floating point operations.}
  The methods should be plotted against both the number of matrix-vector products and an operation estimate.  Under the assumption that multiplication by $A$ or $A^*$ costs $9m$ flops, the main per-step costs are approximately:
  \begin{center}
    \begin{tabular}{lll}
      \toprule
      Method & Matrix-vector products per step & Qualitative extra cost\\
      \midrule
      GMRES & $1$ with $A$ & orthogonalization grows with $n$\\
      GMRES(10) & $1$ with $A$ & orthogonalization bounded by restart length\\
      CGS & $2$ with $A$ & short recurrences\\
      BCGSTAB & $2$ with $A$ & short recurrences and stabilization\\
      QMR & $1$ with $A$ and $1$ with $A^*$ & bi-Lanczos short recurrences\\
      CGNE & $1$ with $A$ and $1$ with $A^*$ & CG on $AA^*$ or $A^*A$\\
      \bottomrule
    \end{tabular}
  \end{center}
  The expected conclusion is:
  \begin{itemize}
    \item Full GMRES normally gives the most reliable residual decrease per iteration, but its orthogonalization cost and storage grow.
    \item GMRES(10) has bounded storage and work per cycle, but restarting can slow or stagnate convergence.
    \item CGS and BCGSTAB are cheap short-recurrence methods.  CGS may be erratic; BCGSTAB is usually more stable.
    \item QMR reduces the irregularity of BCG but requires access to $A^*$.
    \item CGNE applies CG to the normal equations, so the effective condition number is squared; it is often robust but can be slow.
  \end{itemize}
\end{solution}

\begin{problem}\label{pro:5}
  Let $M=LU$ be a preconditioner for $A$.  Show that the left, right, and split preconditioned matrices have the same eigenvalues.  Does this imply the corresponding iterations converge in the same or roughly the same number of steps?
\end{problem}

\begin{solution}
  The three common preconditioned matrices are
  \[
    A_L=M^{-1}A,
    \qquad
    A_R=AM^{-1},
    \qquad
    A_S=L^{-1}AU^{-1}.
  \]
  First, $M^{-1}A$ and $AM^{-1}$ are similar whenever $M$ is nonsingular:
  \[
    AM^{-1}=M(M^{-1}A)M^{-1}.
  \]
  Hence they have exactly the same eigenvalues.

  For the split form, since $M=LU$, we have
  \[
    M^{-1}A=U^{-1}L^{-1}A.
  \]
  The split matrix is
  \[
    L^{-1}AU^{-1}.
  \]
  These two matrices are similar:
  \[
    L^{-1}AU^{-1}=U(U^{-1}L^{-1}A)U^{-1}=U(M^{-1}A)U^{-1}.
  \]
  Therefore
  \[
    \sigma(M^{-1}A)=\sigma(AM^{-1})=\sigma(L^{-1}AU^{-1}).
  \]

  However, equal eigenvalues do not imply identical convergence behavior for nonnormal matrices.  Krylov convergence depends not only on eigenvalues, but also on eigenvectors, field of values, conditioning of the eigenvector matrix, and the residual norm in which optimality is measured.

  Thus:
  \begin{enumerate}[label=(\alph*)]
    \item They do \emph{not} necessarily converge in exactly the same number of steps in finite precision.  In exact arithmetic, unrestarted Krylov methods may terminate in the same degree if the minimal polynomial is the same, but practical convergence histories can differ.
    \item They do \emph{not} necessarily converge in roughly the same number of steps for arbitrary matrices, especially for highly nonnormal matrices.
    \item The most accurate statement is that they may behave similarly when the transformations are well conditioned and the matrix is not strongly nonnormal.  For ill-conditioned or highly nonnormal problems, convergence can be substantially different.
  \end{enumerate}
\end{solution}

\begin{problem}\label{pro:6}
  Let $A$ and its preconditioner $M$ be Hermitian positive definite.  Since $M^{-1}A$ is self-adjoint with respect to the $A$-inner product, write a left-preconditioned CG-like algorithm for
  \[
    M^{-1}Ax=M^{-1}b
  \]
  using the $A$-inner product.  The algorithm should use only one matrix-vector product with $A$ per CG step.
\end{problem}

\begin{solution}
  Define
  \[
    \inner{u}{v}_A=\inner{Au}{v}.
  \]
  For $B=M^{-1}A$, we verify self-adjointness in the $A$-inner product:
  \[
    \inner{Bu}{v}_A
    =\inner{A M^{-1}Au}{v}
    =\inner{u}{A M^{-1}A v}
    =\inner{u}{Bv}_A,
  \]
  because $A$ and $M$ are Hermitian positive definite.  Also
  \[
    \inner{Bu}{u}_A=\inner{A M^{-1}Au}{u}>0
  \]
  for $u\ne0$.  Therefore CG can be applied to $Bx=M^{-1}b$ in the $A$-inner product.

  A convenient implementation is the following.  Let
  \[
    r_k=b-Ax_k,
    \qquad
    z_k=M^{-1}r_k.
  \]
  The preconditioned residual for $M^{-1}Ax=M^{-1}b$ is $z_k$.  The algorithm is:

  \begin{algorithm}[H]
    \caption{$A$-inner-product left-preconditioned CG}
    \begin{algorithmic}[1]
      \State Choose $x_0$; compute $r_0=b-Ax_0$ and solve $Mz_0=r_0$.
      \State Set $p_0=z_0$.
      \For{$k=0,1,2,\ldots$ until convergence}
      \State Compute $w_k=Ap_k$. \Comment{the only multiplication by $A$ in this step}
      \State $\displaystyle \alpha_k=\frac{\inner{r_k}{z_k}}{\inner{p_k}{w_k}}$.
      \State $x_{k+1}=x_k+\alpha_kp_k$.
      \State $r_{k+1}=r_k-\alpha_kw_k$.
      \State Solve $Mz_{k+1}=r_{k+1}$.
      \State $\displaystyle \beta_k=\frac{\inner{r_{k+1}}{z_{k+1}}}{\inner{r_k}{z_k}}$.
      \State $p_{k+1}=z_{k+1}+\beta_kp_k$.
      \EndFor
    \end{algorithmic}
  \end{algorithm}

  Although this looks like standard preconditioned CG, its derivation is from applying CG to $M^{-1}A$ under the $A$-inner product.  The formula uses only one $A$-multiplication per iteration because $Ap_k$ is reused for both the step length and the residual update.
\end{solution}

\begin{problem}\label{pro:7}
  Describe right and split preconditioned CGNR/CGNE, and the two centered normal-equation versions.
\end{problem}

\begin{solution}
  Recall the two normal-equation formulations:
  \[
    \text{CGNR: } A^*Ax=A^*b,
    \qquad
    \text{CGNE: } AA^*y=b,
    \quad x=A^*y.
  \]

  \medskip
  \noindent\textbf{(a) Four right/split preconditioned versions.}

  \textbf{Right P-CGNR.}
  Let $x=M^{-1}u$.  Then
  \[
    A M^{-1}u=b.
  \]
  Applying CGNR gives
  \[
    (AM^{-1})^*(AM^{-1})u=(AM^{-1})^*b,
  \]
  that is
  \[
    M^{-*}A^*AM^{-1}u=M^{-*}A^*b,
    \qquad x=M^{-1}u.
  \]

  \textbf{Right P-CGNE.}
  Again set $x=M^{-1}u$.  CGNE for $AM^{-1}u=b$ is
  \[
    (AM^{-1})(AM^{-1})^*y=b,
    \qquad u=(AM^{-1})^*y.
  \]
  Thus
  \[
    AM^{-1}M^{-*}A^*y=b,
    \qquad
    x=M^{-1}u=M^{-1}M^{-*}A^*y.
  \]

  \textbf{Split P-CGNR.}
  Let $M=LU$ and set $x=U^{-1}u$.  The split-preconditioned system is
  \[
    L^{-1}AU^{-1}u=L^{-1}b.
  \]
  CGNR gives
  \[
    (L^{-1}AU^{-1})^*(L^{-1}AU^{-1})u
    =(L^{-1}AU^{-1})^*L^{-1}b.
  \]
  Equivalently,
  \[
    U^{-*}A^*L^{-*}L^{-1}AU^{-1}u
    =U^{-*}A^*L^{-*}L^{-1}b,
    \qquad x=U^{-1}u.
  \]

  \textbf{Split P-CGNE.}
  For
  \[
    L^{-1}AU^{-1}u=L^{-1}b,
  \]
  CGNE gives
  \[
    (L^{-1}AU^{-1})(L^{-1}AU^{-1})^*y=L^{-1}b,
    \qquad
    u=(L^{-1}AU^{-1})^*y.
  \]
  Thus
  \[
    L^{-1}AU^{-1}U^{-*}A^*L^{-*}y=L^{-1}b,
  \]
  and
  \[
    x=U^{-1}u
    =U^{-1}U^{-*}A^*L^{-*}y.
  \]
  Suitable weighted inner products can be used so that the coefficient matrices are explicitly Hermitian positive definite in the chosen inner product.

  \medskip
  \noindent\textbf{(b) Centered versions.}

  \textbf{Centered NE.}
  The system is
  \[
    AM^{-1}A^*y=b,
    \qquad
    x=M^{-1}A^*y.
  \]
  Since $M$ is Hermitian positive definite, $AM^{-1}A^*$ is Hermitian positive semidefinite, and positive definite if $A$ is nonsingular.  CG can be applied directly:
  \begin{algorithm}[H]
    \caption{Centered NE-CG}
    \begin{algorithmic}[1]
      \State Choose $y_0$; set $x_0=M^{-1}A^*y_0$ and $r_0=b-AM^{-1}A^*y_0$.
      \State $p_0=r_0$.
      \For{$k=0,1,2,\ldots$}
      \State $q_k=A M^{-1}A^*p_k$.
      \State $\alpha_k=(r_k^*r_k)/(p_k^*q_k)$.
      \State $y_{k+1}=y_k+\alpha_kp_k$.
      \State $r_{k+1}=r_k-\alpha_kq_k$.
      \State $\beta_k=(r_{k+1}^*r_{k+1})/(r_k^*r_k)$.
      \State $p_{k+1}=r_{k+1}+\beta_kp_k$.
      \EndFor
      \State Return $x=M^{-1}A^*y$.
    \end{algorithmic}
  \end{algorithm}

  \textbf{Centered NR.}
  The system is
  \[
    A^*M^{-1}Ax=A^*M^{-1}b.
  \]
  Apply CG directly to this Hermitian positive definite system:
  \begin{algorithm}[H]
    \caption{Centered NR-CG}
    \begin{algorithmic}[1]
      \State Choose $x_0$; set $r_0=A^*M^{-1}(b-Ax_0)$.
      \State $p_0=r_0$.
      \For{$k=0,1,2,\ldots$}
      \State $q_k=A^*M^{-1}Ap_k$.
      \State $\alpha_k=(r_k^*r_k)/(p_k^*q_k)$.
      \State $x_{k+1}=x_k+\alpha_kp_k$.
      \State $r_{k+1}=r_k-\alpha_kq_k$.
      \State $\beta_k=(r_{k+1}^*r_{k+1})/(r_k^*r_k)$.
      \State $p_{k+1}=r_{k+1}+\beta_kp_k$.
      \EndFor
    \end{algorithmic}
  \end{algorithm}
\end{solution}

\begin{problem}\label{pro:8}
  Assume a Hermitian positive definite matrix $M$ is used to precondition GMRES for a non-Hermitian system.  Give a formal description of left-preconditioned GMRES with the $M$-inner product, including modified Gram-Schmidt.  What optimality property holds?  What if $A$ is also Hermitian?  What is a potential advantage?
\end{problem}

\begin{solution}
  Consider the left-preconditioned system
  \[
    M^{-1}Ax=M^{-1}b.
  \]
  Define the $M$-inner product and norm by
  \[
    \inner{u}{v}_M=v^*Mu,
    \qquad
    \norm{u}_M=\sqrt{\inner{u}{u}_M}.
  \]
  Let
  \[
    B=M^{-1}A,
    \qquad
    c=M^{-1}b.
  \]
  GMRES with the $M$-inner product constructs an $M$-orthonormal basis
  \[
    Q_n=[q_1,\ldots,q_n],
    \qquad
    Q_n^*MQ_n=I.
  \]
  Because inner products require $Mq_i$, we store both $q_i$ and $z_i=Mq_i$.

  \begin{algorithm}[H]
    \caption{Left-preconditioned GMRES with $M$-inner product}
    \begin{algorithmic}[1]
      \State Choose $x_0$ and compute $r_0=c-Bx_0=M^{-1}(b-Ax_0)$.
      \State Compute $z=M r_0=b-Ax_0$.
      \State $\beta=\norm{r_0}_M=\sqrt{r_0^*Mr_0}=\sqrt{r_0^*(b-Ax_0)}$.
      \State $q_1=r_0/\beta$ and $z_1=Mq_1=z/\beta$.
      \For{$j=1,2,\ldots,n$}
      \State Compute $w=Bq_j=M^{-1}Aq_j$ by first forming $Aq_j$ and then solving with $M$.
      \State Store $s=Mw=Aq_j$.
      \For{$i=1,2,\ldots,j$} \Comment{modified Gram-Schmidt in $M$-inner product}
      \State $h_{ij}=q_i^*Mw=q_i^*s=z_i^*w$.
      \State $w=w-h_{ij}q_i$.
      \State $s=s-h_{ij}z_i$. \Comment{maintain $s=Mw$}
      \EndFor
      \State $h_{j+1,j}=\sqrt{w^*s}$.
      \State $q_{j+1}=w/h_{j+1,j}$, \quad $z_{j+1}=s/h_{j+1,j}$.
      \EndFor
      \State Solve the least-squares problem
      \[
        y_n=\arg\min_y\norm{\beta e_1-\bar H_ny}_2.
      \]
      \State Return $x_n=x_0+Q_ny_n$.
    \end{algorithmic}
  \end{algorithm}

  The Arnoldi relation is
  \[
    BQ_n=Q_{n+1}\bar H_n,
    \qquad
    Q_{n+1}^*MQ_{n+1}=I.
  \]
  Therefore
  \begin{align*}
    \norm{c-B(x_0+Q_ny)}_M
    &=\norm{r_0-BQ_ny}_M\\
    &=\norm{Q_{n+1}(\beta e_1-\bar H_ny)}_M\\
    &=\norm{\beta e_1-\bar H_ny}_2.
  \end{align*}
  Thus the computed $x_n$ satisfies the optimality property
  \[
    \boxed{
      x_n=\arg\min_{x\in x_0+\K_n(M^{-1}A,r_0)}
      \norm{M^{-1}(b-Ax)}_M
    }
  \]
  or equivalently
  \[
    \norm{M^{-1}(b-Ax)}_M^2
    =(b-Ax)^*M^{-1}(b-Ax).
  \]

  If $A$ is also Hermitian, then $M^{-1}A$ is self-adjoint with respect to the $M$-inner product:
  \[
    \inner{M^{-1}Au}{v}_M
    =v^*A u
    =\inner{u}{M^{-1}Av}_M.
  \]
  Hence the Arnoldi process reduces to a three-term Lanczos process in the $M$-inner product.  The resulting method is therefore a minimum-residual method for a self-adjoint preconditioned operator.  A potential advantage is that one can exploit short recurrences, reducing storage and orthogonalization cost compared with full GMRES, while still measuring residuals in a norm naturally induced by the preconditioner.
\end{solution}

\end{document}
